For BERT model, we first performed hyperparameter tuning on the training set. BERT was tuned using grid search over the parameter ranges shown in Table~\ref{tab:hyperparameters}. Training each configuration with early stopping based on validation loss, we selected the best hyperparameters based on macro-F1 on the development set. 

% During tuning, each configuration was trained for 5 epochs, and the best configuration was selected based on macro-F1 on the development set. Early stopping with a patience of 2 epochs was enabled during this tuning stage, although in practice the short tuning runs meant that most configurations finished all 5 epochs without triggering early stopping.

After selecting the best hyperparameters for a given model, we trained that model again on the train data and evaluated it on the test set. The final training runs used the Adam optimiser, a batch size of 512, gradient clipping with a maximum norm of 1.0, and a maximum of 15 epochs. Unlike the tuning stage, the final training procedure did not apply early stopping; instead, models were trained for the full configured number of epochs and then evaluated on the test set. We made this choice to obtain directly comparable learning curves across all runs and to analyse convergence behaviour in more detail. However, this design introduces a tradeoff: by prioritising diagnostic information, we may give up some generalisation performance relative to selecting the best validation checkpoint with early stopping. The final test scores should therefore be interpreted as results from a fixed-length diagnostic training protocol, not necessarily as the best achievable numbers for each architecture.

We report accuracy and macro-F1 as our main evaluation metrics, with macro-F1 used as the primary measure for model selection because it gives equal weight to all four classes. In addition, we generated classification reports, confusion matrices, learning curves, and files with misclassified test examples for later error analysis. The test set was only used for final evaluation after model selection.

As our ablation study, we varied the maximum input sequence length while keeping the rest of the hyperparameters fixed to the respective base models' constants (same setup). We ran each model with maximum sequence lengths of 64, 128, and 256 tokens, resulting in six main experiment settings: CNN-seq64, CNN-seq128, CNN-seq256, LSTM-seq64, LSTM-seq128, and LSTM-seq256. This design allows us to compare the two architectures directly and also study how much performance depends on the amount of input context available to the model.
